\subsection{When a function can be reversed}

An inverse function reverses an assignment. If \(f\) sends \(x\) to \(y\), then \(f^{-1}\) should send \(y\) back to \(x\). This is possible only when different allowed inputs do not produce the same output.

\begin{definitionbox}[One-to-one function]
A function is one-to-one on its domain if
\[
f(x_1)=f(x_2)
\quad\Longrightarrow\quad
x_1=x_2.
\]
Only a one-to-one function can have an inverse on its entire range.
\end{definitionbox}

\begin{examplebox}[Why \(x^2\) needs a restricted domain]
On all real numbers, \(f(x)=x^2\) is not one-to-one because \(f(2)=f(-2)=4\). On the restricted domain \([0,\infty)\), however, every output has one non-negative input, so the inverse is \(f^{-1}(y)=\sqrt y\).
\end{examplebox}

\begin{interpretationbox}[Horizontal-line test]
A graph represents a one-to-one function precisely when every horizontal line meets it at most once. The test asks whether an output identifies a unique input.
\end{interpretationbox}
